%% 中文摘要页


\begin{center}\heiti\sanhao\textbf{
        摘要
    }
\end{center}

随着无人化装备向低成本、智能化方向发展，在移动平台上实现对固定目标的实时识别与精确瞄准已成为领域研究热点。现有方案普遍存在成本高昂、部署复杂或实时性不足的问题。本文提出并实现了一套基于移动平台的固定目标识别及瞄准系统，以``千元级成本、纯视觉+IMU 紧耦合''为核心约束，设计了可工程化落地的低成本解决方案。

系统采用三节点分布式异构架构：Linux 视觉节点负责目标检测、PnP 位姿解算与跟踪状态管理；STM32H743 云台主控节点以 480\,MHz 主频运行 RT-Thread 实时操作系统，执行双闭环姿态控制与 IMU 前馈补偿；MSPM0G3507 底盘主控节点独立完成八路灰度循迹与差速编码器 PI 速度闭环，三节点通过串口协议协同工作，互不干扰。

在视觉识别与跟踪算法上，本文采用 HSV 颜色阈值分割、轮廓几何约束与 PnP 重投影尺寸核查三级筛选策略，对 A4 白纸目标实现鲁棒检测与中心定位，并以三线程+双有界队列结构保证 25 FPS 以上的视觉吞吐稳定性。

在云台控制算法上，本文提出``像素误差外环 APID + IMU 角速度/加速度前馈内环''的双闭环结构，前馈项参考自抗扰控制思想，在平台突加速或振动工况下可显著降低相位滞后，实测可将瞄准线峰值偏差压缩 30\% 以上，稳态指向误差满足 $\leq$0.8\,mrad 的设计指标。IMU 采用 HiPNUC 六轴惯性传感器，以 100\,Hz 输出角速度与姿态数据，通过 \texttt{hipnuc\_dec} 协议模块完成字节流解析与 CRC 校验。

在底盘控制上，系统采用差速运动学模型完成速度分解，左右轮分别执行增量式 PI 调速与 AB 相编码器反馈闭环；循迹策略融合归一化灰度加权偏差估计与一阶滞后低通滤波，有效抑制传感噪声引起的控制抖动。此外，系统集成了 A* 启发式路径规划算法，可在结构化场地中完成点对点自主机动。

实验结果表明，系统在室内标准光照条件下视觉帧率稳定超过 25 FPS，IMU 采样频率维持 100\,Hz，底盘可稳定完成多圈循迹任务，整机器件成本控制在千元级别，验证了所提方案的可行性与工程实用价值。

\textbf{关键词：}视觉伺服，目标识别，云台控制，IMU 前馈，移动机器人

\clearpage